\documentclass[a4paper]{article}
\usepackage{iftex}
\ifPDFTeX
  \usepackage{vntex}
\else
  % XeLaTeX / LuaLaTeX (vi du Overleaf chon compiler khac): dung fontspec
  \usepackage{fontspec}
  \setmainfont{TeX Gyre Termes}
  \setsansfont{TeX Gyre Heros}
  \setmonofont{TeX Gyre Cursor}
\fi
%\usepackage[english,vietnam]{babel}
%\usepackage[utf8]{inputenc}

%\usepackage[utf8]{inputenc}
%\usepackage[francais]{babel}
\usepackage{a4wide,amssymb,epsfig,latexsym,multicol,array,hhline,fancyhdr}
\usepackage{booktabs}
\usepackage{longtable}
\usepackage{amsmath}
\usepackage{lastpage}
\usepackage[lined,boxed,commentsnumbered]{algorithm2e}
\usepackage{enumerate}
\usepackage{color}
\usepackage{graphicx}							% Standard graphics package
\usepackage{array}
\usepackage{tabularx, caption}
\usepackage{multirow}
\usepackage[framemethod=tikz]{mdframed}% For highlighting paragraph backgrounds
\usepackage{multicol}
\usepackage{rotating}
\usepackage{graphics}
\usepackage{geometry}
\usepackage{setspace}
\usepackage{epsfig}
\usepackage{tikz}
\usepackage{listings}
\usetikzlibrary{arrows,snakes,backgrounds}
\usepackage{hyperref}
\hypersetup{urlcolor=blue,linkcolor=black,citecolor=black,colorlinks=true} 
%\usepackage{pstcol} 								% PSTricks with the standard color package

\newtheorem{theorem}{{\bf Định lý}}
\newtheorem{property}{{\bf Tính chất}}
\newtheorem{proposition}{{\bf Mệnh đề}}
\newtheorem{corollary}[proposition]{{\bf Hệ quả}}
\newtheorem{lemma}[proposition]{{\bf Bổ đề}}

\everymath{\color{blue}}
%\usepackage{fancyhdr}
\setlength{\headheight}{40pt}
\pagestyle{fancy}
\fancyhead{} % clear all header fields
\fancyhead[L]{
 \begin{tabular}{rl}
    \begin{picture}(25,15)(0,0)
    \put(0,-8){\includegraphics[width=8mm, height=8mm]{logoITSGUsmall.png}}
    %\put(0,-8){\epsfig{width=10mm,figure=hcmut.eps}}
   \end{picture}&
	%\includegraphics[width=8mm, height=8mm]{hcmut.png} & %
	\begin{tabular}{l}
		\textbf{\bf \ttfamily Trường Đại học Sài Gòn}\\
		\textbf{\bf \ttfamily Khoa Công Nghệ Thông Tin}
	\end{tabular} 	
 \end{tabular}
}
\fancyhead[R]{
	\begin{tabular}{l}
		\tiny \bf \\
		\tiny \bf 
	\end{tabular}  }
\fancyfoot{} % clear all footer fields
\fancyfoot[L]{\scriptsize \ttfamily Bài tập lớn 2 (cuối kì) môn Công nghệ phần mềm - Năm học 2026-2027}
\fancyfoot[R]{\scriptsize \ttfamily Trang {\thepage}/\pageref{LastPage}}
\renewcommand{\headrulewidth}{0.3pt}
\renewcommand{\footrulewidth}{0.3pt}


%%%
\setcounter{secnumdepth}{4}
\setcounter{tocdepth}{3}
\makeatletter
\newcounter {subsubsubsection}[subsubsection]
\renewcommand\thesubsubsubsection{\thesubsubsection .\@alph\c@subsubsubsection}
\newcommand\subsubsubsection{\@startsection{subsubsubsection}{4}{\z@}%
                                     {-3.25ex\@plus -1ex \@minus -.2ex}%
                                     {1.5ex \@plus .2ex}%
                                     {\normalfont\normalsize\bfseries}}
\newcommand*\l@subsubsubsection{\@dottedtocline{3}{10.0em}{4.1em}}
\newcommand*{\subsubsubsectionmark}[1]{}
\makeatother

\definecolor{dkgreen}{rgb}{0,0.6,0}
\definecolor{gray}{rgb}{0.5,0.5,0.5}
\definecolor{mauve}{rgb}{0.58,0,0.82}

\lstset{frame=tb,
	language=Matlab,
	aboveskip=3mm,
	belowskip=3mm,
	showstringspaces=false,
	columns=flexible,
	basicstyle={\small\ttfamily},
	numbers=none,
	numberstyle=\tiny\color{gray},
	keywordstyle=\color{blue},
	commentstyle=\color{dkgreen},
	stringstyle=\color{mauve},
	breaklines=true,
	breakatwhitespace=true,
	tabsize=3,
	numbers=left,
	stepnumber=1,
	numbersep=1pt,    
	firstnumber=1,
	numberfirstline=true
}

\newcommand{\fillme}[1]{{\color{red}[#1]}}
\newcommand{\hinhx}[3]{\begin{center}\includegraphics[width=#2,height=0.66\textheight,keepaspectratio]{#1}\\[2pt]{\small\textit{#3}}\end{center}}
\newcommand{\hinh}[4]{\begin{figure}[htbp]\centering\includegraphics[width=#2,height=0.6\textheight,keepaspectratio]{#1}\caption{#3}\label{#4}\end{figure}}
\newcommand{\hinhngang}[3]{\begin{sidewaysfigure}[p]\centering\includegraphics[width=0.97\textheight,height=0.78\textwidth,keepaspectratio]{#1}\caption{#2}\label{#3}\end{sidewaysfigure}}
\newcommand{\hinhcao}[3]{\begin{figure}[htbp]\centering\includegraphics[height=0.78\textheight,width=\linewidth,keepaspectratio]{#1}\caption{#2}\label{#3}\end{figure}}
\newcommand{\nguon}[1]{\par\vspace{-4pt}{\centering\footnotesize\textit{Nguồn: #1}\par}\medskip}
\graphicspath{{figures/}{uml/}}
\everymath{}
\begin{document}

\begin{titlepage}
\begin{center}
TRƯỜNG ĐẠI HỌC SÀI GÒN \\
KHOA CÔNG NGHỆ THÔNG TIN
\end{center}
\vspace{1cm}

\begin{figure}[h!]
\begin{center}
\includegraphics[width=3cm]{logoITSGU.png}
\end{center}
\end{figure}

\vspace{1cm}

\begin{center}
\begin{tabular}{c}
	\multicolumn{1}{l}{\textbf{{\Large CÔNG NGHỆ PHẦN MỀM}}}\\
	~~\\
	\hline
	\\
	\multicolumn{1}{l}{\textbf{{\large Bài tập lớn 2 (cuối kì): từ yêu cầu đến triển khai}}}\\
	\\
	\textbf{{\huge BrewLite -- Ứng dụng đặt đồ uống}}\\
	\\
	\hline
\end{tabular}
\end{center}

\vspace{3cm}

\begin{table}[h]
\begin{tabular}{rrl}
\hspace{5 cm} & GVHD: &Từ Lăng Phiêu\\
& SV: & Nguyễn Lê Nhựt Thắng - 3124560085\\
& & Hồ Thiên Ngân - 3124410224 \\
& & Hà Thiên Long - 3124410193 \\
& & Giản Thiên Lực - 3124410206\\
\end{tabular}
\vspace{1.5 cm}
\end{table}

\begin{center}
{\footnotesize TP. HỒ CHÍ MINH, THÁNG 11/2026}
\end{center}
\end{titlepage}

\thispagestyle{empty}

\newpage
\tableofcontents
\clearpage
\listoffigures
\clearpage
\listoftables
\clearpage
\section*{Danh mục từ viết tắt}\addcontentsline{toc}{section}{Danh mục từ viết tắt}
\begin{longtable}{@{}p{3cm}p{10cm}@{}}
\toprule
\textbf{Viết tắt} & \textbf{Nghĩa đầy đủ} \\ \midrule \endhead
ADR & Architecture Decision Record (bản ghi quyết định kiến trúc) \\
API & Application Programming Interface \\
CD & Continuous Delivery / Deployment (phân phối, triển khai liên tục) \\
CI & Continuous Integration (tích hợp liên tục) \\
CSDL & Cơ sở dữ liệu \\
DoD & Definition of Done (định nghĩa hoàn thành) \\
E2E & End-to-End (kiểm thử đầu cuối) \\
ERD & Entity Relationship Diagram (sơ đồ thực thể -- quan hệ) \\
FR / NFR & Functional / Non-functional Requirement (yêu cầu chức năng / phi chức năng) \\
GHCR & GitHub Container Registry \\
JWT & JSON Web Token \\
MVP & Minimum Viable Product (sản phẩm khả dụng tối thiểu) \\
PO & Product Owner \\
PR & Pull Request \\
SRS & Software Requirements Specification (đặc tả yêu cầu phần mềm) \\
UC & Use Case \\
UI & User Interface (giao diện người dùng) \\
UML & Unified Modeling Language \\
\bottomrule
\end{longtable}

\clearpage
\section*{Bảng thuật ngữ}\addcontentsline{toc}{section}{Bảng thuật ngữ}
\begin{longtable}{@{}p{3.3cm}p{3.3cm}p{6.4cm}@{}}
\toprule
\textbf{Thuật ngữ (Việt)} & \textbf{Tiếng Anh} & \textbf{Định nghĩa ngắn} \\ \midrule \endhead
Idempotency-Key & Idempotency key & Khóa do client gửi kèm yêu cầu để gửi lại cùng yêu cầu không tạo kết quả trùng. \\
Quay lui & Rollback & Triển khai lại một phiên bản ổn định trước đó khi bản mới có sự cố. \\
Kiểm thử khói & Smoke test & Kiểm tra nhanh sau triển khai rằng dịch vụ còn sống (ở đây là gọi \texttt{/health}). \\
Máy trạng thái & State machine & Tập trạng thái của đơn hàng và các chuyển trạng thái hợp lệ. \\
Nghiệm thu chéo & Cross acceptance & Người khác với người viết task kiểm tra task đạt DoD. \\
\bottomrule
\end{longtable}

\clearpage
\section*{Tóm tắt}\addcontentsline{toc}{section}{Tóm tắt}
Báo cáo trình bày Assignment 2 của dự án BrewLite, đi từ yêu cầu đến triển khai: khai thác và đặc tả yêu cầu (17 yêu cầu chức năng, 11 yêu cầu phi chức năng), mô hình hóa bằng UML, thiết kế kiến trúc kèm bảy quyết định kiến trúc, và hai Sprint hiện thực (20--26/10/2026). Sản phẩm là ứng dụng đặt đồ uống dùng Next.js, NestJS, PostgreSQL, có pipeline CI bốn job, pipeline CD từ tag phiên bản lên GHCR và Render, kèm smoke test. Nhóm đã chứng minh CI chặn merge khi test lỗi và quay lui thành công về bản trước rồi triển khai lại. Phiên bản bàn giao là \texttt{v0.1.3} (07/11/2026). Ma trận truy vết nối yêu cầu với thiết kế, mã nguồn và test. Còn mở: dấu vết nghiệm thu chéo trên GitHub, và các giới hạn của gói hạ tầng miễn phí (web tạm ngủ, CSDL hết hạn 22/11/2026).

\vspace{6pt}
\noindent\textbf{Từ khóa:} yêu cầu phần mềm; UML; kiến trúc; CI/CD; triển khai; truy vết

\clearpage


%%%%%%%%%%%%%%%%%%%%%%%%%%%%%%%%%
\section{Giới thiệu và quy trình phát triển}\label{ch:gioithieu}
Bảng~\ref{tab:doichieu-rutgon} cho biết yêu cầu của Assignment 2 nằm ở chương nào; bản đầy đủ ở mục~\ref{app:doichieu}.
\begin{table}[htbp]\centering
\caption{Đối chiếu rút gọn giữa Task của đề và chương báo cáo}\label{tab:doichieu-rutgon}
\begin{tabular}{@{}p{6.4cm}p{6.6cm}@{}}\toprule
\textbf{Task của đề} & \textbf{Chương} \\ \midrule
Task 1: khai thác và đặc tả yêu cầu & mục~\ref{ch:yeucau} \\
Task 2: mô hình hóa (activity, sequence, class) & mục~\ref{ch:mohinh} \\
Task 3: kiến trúc và triển khai (deployment) & mục~\ref{ch:kientruc} \\
Task 4: Sprint 1, repository, MVP, CI & mục~\ref{ch:quanly}, mục~\ref{ch:trienkhai-se} \\
Task 5: Sprint 2, CI/CD, bàn giao & mục~\ref{ch:trienkhai-se}, mục~\ref{ch:kiemthu}, mục~\ref{ch:banggiao} \\
Truy vết yêu cầu & mục~\ref{ch:truyvet} \\ \bottomrule
\end{tabular}
\end{table}

\subsection{Bối cảnh, mục tiêu, phạm vi}\label{sec:se-boicanh}
Quán cà phê quy mô nhỏ cần để khách đặt món và thanh toán không dùng tiền mặt ngay trên điện thoại, tránh xếp hàng và nhầm đơn. BrewLite là ứng dụng web mobile-first cho phép xem menu, chọn size và topping, quản lý giỏ hàng, đăng nhập, đặt đơn, thanh toán mô phỏng và xem lịch sử đơn.

Assignment 2 tiếp nối Assignment 1 và đi xuyên suốt từ khai thác yêu cầu, mô hình hóa, thiết kế kiến trúc, hiện thực hai Sprint, CI/CD đến bàn giao bản production. Mã nguồn dùng chung với Assignment 1; báo cáo này tập trung vào đặc tả, thiết kế, quy trình và vận hành. Phạm vi chi tiết (trong, ngoài, không thực hiện) ở mục~\ref{sec:se-phamvi}.

\subsection{Mô hình quy trình áp dụng}\label{sec:se-quytrinh}
Nhóm áp dụng Agile Scrum~\cite{scrum}: Product Backlog, Sprint Backlog, Increment; bốn sự kiện Sprint Planning, Daily Scrum, Sprint Review, Sprint Retrospective. Scrum phù hợp vì thời gian ngắn (khoảng một tuần), yêu cầu có thể điều chỉnh sau mỗi Sprint và cần sản phẩm chạy được sớm để triển khai thử. Assignment 2 chia thành hai Sprint: Sprint 1 (20--23/10) tạo nền tảng và MVP menu, Sprint 2 (24--26/10) hoàn thiện luồng end-to-end, CI/CD và bàn giao (xem mục~\ref{sec:se-sprint}).

\subsection{Công nghệ và công cụ sử dụng}\label{sec:se-congnghe}
\begin{table}[htbp]\centering
\caption{Công nghệ và công cụ}\label{tab:congnghe}
\begin{tabular}{@{}p{3.2cm}p{9.6cm}@{}}\toprule
\textbf{Lớp} & \textbf{Công nghệ} \\ \midrule
Frontend & Next.js (App Router), mobile-first, manifest PWA cơ bản \\
Backend & NestJS, class-validator, JWT, bcrypt \\
Cơ sở dữ liệu & PostgreSQL 16, Prisma 7 (adapter \texttt{pg}), migration và seed \\
Kiểm thử & Vitest, Supertest (PostgreSQL thật), Playwright \\
Vận hành & Docker Compose, GitHub Actions, GHCR, Render \\
Sơ đồ & PlantUML (nguồn trong \path{docs/diagrams}) \\ \bottomrule
\end{tabular}
\end{table}
\begin{itemize}\setlength{\itemsep}{0pt}
\item Kho mã nguồn: \url{https://github.com/tthan-prog/BrewLite}
\item Web production: \url{https://brewlite-web.onrender.com}
\item API: \url{https://brewlite-api.onrender.com/health}
\item Phiên bản bàn giao: tag \texttt{v0.1.3}, phát hành 07/11/2026\\(commit \texttt{ecd11bd}, image \texttt{sha-ecd11bd})
\end{itemize}

\subsection{Cấu trúc báo cáo}\label{sec:se-cautruc}
mục~\ref{ch:quanly} trình bày quản lý dự án; mục~\ref{ch:yeucau} đặc tả yêu cầu; mục~\ref{ch:mohinh} mô hình UML; mục~\ref{ch:kientruc} kiến trúc; mục~\ref{ch:trienkhai-se} hiện thực và CI/CD; mục~\ref{ch:kiemthu} kiểm thử; mục~\ref{ch:banggiao} bàn giao và nhìn lại; mục~\ref{ch:truyvet} ma trận truy vết; mục~\ref{ch:ketluan-se} kết luận. Phụ lục gồm bảng đối chiếu đề bài, phân công và đóng góp, và tư liệu bổ sung.

\section{Quản lý dự án (Scrum)}\label{ch:quanly}
\subsection{Vai trò / trách nhiệm trong nhóm}\label{sec:se-vaitro}
Vai trò theo \path{docs/chung/team-agreement.md}. Đây là ``người chịu trách nhiệm chính'', không phải chỉ một người làm. Scrum Guide 2020 gọi các vai là \emph{accountabilities}; báo cáo dùng từ ``vai trò'' theo đề bài.
\begin{table}[htbp]\centering
\caption{Vai trò và trách nhiệm}\label{tab:vaitro}
\begin{tabular}{@{}p{3.4cm}p{3.8cm}p{5.8cm}@{}}\toprule
\textbf{Thành viên} & \textbf{Vai trò} & \textbf{Trách nhiệm cụ thể} \\ \midrule
Nguyễn Lê Nhựt Thắng (A) & Scrum Master, phát triển backend và nền tảng & Repo, CI/CD, triển khai, điều phối Scrum; API sản phẩm, xác thực, thanh toán; sequence diagram \\
Hồ Thiên Ngân (B) & Backend & Dữ liệu, seed, tạo đơn, trạng thái đơn \\
Hà Thiên Long (C) & Kiểm thử, Product Owner tạm thời & Tài liệu yêu cầu, truy vết, backlog; unit, tích hợp, E2E \\
Giản Thiên Lực (D) & Frontend & Menu, giỏ, thanh toán, lịch sử; activity, use-case; slide, video \\
\bottomrule
\end{tabular}
\end{table}

\subsection{Sự kiện và artefact}\label{sec:se-suKien}
Sprint Planning đầu Sprint (20/10 và 24/10), Daily Scrum hằng ngày tối đa 15 phút qua nhóm chat, Sprint Review và Retrospective cuối Sprint (24/10 và 26/10). Artefact: Product Backlog, Sprint Backlog (kèm Sprint Goal) và Increment là bản chạy được trên production. Product Backlog và Sprint Backlog tổng hợp ở \path{docs/as2/scrum/product-backlog.md}; biên bản nằm ở \path{docs/as2/scrum/minutes/a2-*}. Quy ước chung: không đẩy trực tiếp lên \texttt{main}; mỗi chức năng một nhánh; PR phải qua check \texttt{api} và \texttt{web}, cần một approval.

\subsection{Kế hoạch và diễn biến từng Sprint}\label{sec:se-sprint}
\subsubsection{Sprint 1 (20--23/10/2026)}
\textbf{Sprint Goal:} tạo nền tảng kỹ thuật, tài liệu có kiểm soát phiên bản và MVP đầu tiên cho luồng xem menu. Planning họp ngày 20/10.
\begin{longtable}{@{}p{1cm}p{6.2cm}p{3.2cm}p{2.6cm}@{}}
\caption{Sprint 1: phạm vi và người nghiệm thu}\label{tab:sprint1}\\
\toprule \textbf{Task} & \textbf{Nội dung} & \textbf{Làm chính} & \textbf{Nghiệm thu} \\ \midrule \endfirsthead
\toprule \textbf{Task} & \textbf{Nội dung} & \textbf{Làm chính} & \textbf{Nghiệm thu} \\ \midrule \endhead
\bottomrule \endlastfoot
4.1 & Repo, cấu trúc, README, \texttt{.env.example}, branch protection, quy ước commit, DoD & A & B \\
4.2 & Tạo phẩm Task 1--3 trong \path{docs}, tag \texttt{assign1-final} (commit \texttt{9581384}), thay đổi qua PR & A; C (viết yêu cầu) & D \\
4.3 & Dữ liệu mẫu, \texttt{GET /products}, \texttt{GET /products/:id} & B (dữ liệu), A (API) & C \\
4.3 & Màn hình menu mobile-first, loading, empty, error & D & A \\
4.3 & Kiểm thử tối thiểu API và menu & C & A \\
4.4 & Pipeline CI, chặn merge khi đỏ & A & C \\
\end{longtable}
Kết quả Review (24/10): cả bốn hạng mục demo được và đạt DoD, Sprint Goal đạt. Điểm còn lại cuối Sprint là 0 trên 4 việc; repo chưa có file burndown riêng cho Assignment 2 nên số này tính theo số việc, không phải story point.

\subsubsection{Sprint 2 (24--26/10/2026)}
\textbf{Sprint Goal:} hoàn thành luồng end-to-end, tự động hóa đóng gói và triển khai, có bản production truy cập được. Planning họp ngày 24/10.
\begin{longtable}{@{}p{1cm}p{6.2cm}p{3.2cm}p{2.6cm}@{}}
\caption{Sprint 2: phạm vi và người nghiệm thu}\label{tab:sprint2}\\
\toprule \textbf{Task} & \textbf{Nội dung} & \textbf{Làm chính} & \textbf{Nghiệm thu} \\ \midrule \endfirsthead
\toprule \textbf{Task} & \textbf{Nội dung} & \textbf{Làm chính} & \textbf{Nghiệm thu} \\ \midrule \endhead
\bottomrule \endlastfoot
5.1 & Đăng ký, đăng nhập, thanh toán mô phỏng & A & B \\
5.1 & Tạo đơn, trạng thái đơn & B & C \\
5.1 & Giao diện menu, giỏ, thanh toán, lịch sử & D & A \\
5.2 & Unit, integration, E2E, Docker Compose, \texttt{/health} & C & B \\
5.3 & Pipeline CD, rollback & A & B \\
5.4 & Triển khai production, release \texttt{v0.1.x}, README & A & D \\
5.4 & Video demo end-to-end & D & C \\
5.4 & Sprint Review và Retrospective & A & Cả nhóm \\
\end{longtable}
Kết quả Review (26/10): Task 5.1, 5.2 và 5.3 demo được và đạt DoD; Task 5.4 ``gần đạt''; video demo chưa xong tại thời điểm Review và được hoàn tất sau đó (xem mục~\ref{sec:se-sanpham}). Điểm còn lại lúc Review là 2 việc (video), tính theo số việc.

\subsection{Definition of Done}\label{sec:se-dod}
Bốn mục DoD áp dụng cho mọi task (\path{docs/as1/scrum/dod-nghiem-thu.md}): (D1) chạy được, không lỗi build; (D2) commit rõ ràng và PR có ít nhất một Approve của người khác; (D3) API validate và giao diện có loading, lỗi; (D4) README có hướng dẫn chạy. Người nghiệm thu khác người viết task.

\textbf{Mức độ áp dụng trung thực:} D1, D3, D4 có bằng chứng (CI xanh, test, README). Với D2, nhóm ghi nhận nghiệm thu chéo trong tài liệu, nhưng các PR chưa có Approve hoặc bình luận nghiệm thu của người đánh giá trên GitHub (cột Reviews trống). Điểm này đã được nêu ở Retrospective của cả hai Sprint.

\section{Đặc tả yêu cầu}\label{ch:yeucau}
Đặc tả đầy đủ nằm trong thư mục \path{docs/requirements} (các tệp \path{srs.md}, \path{context-and-scope.md} và \path{use-case-checkout.md}). Hệ mã FR/NFR được giữ xuyên suốt để lập ma trận truy vết ở mục~\ref{ch:truyvet}.

\subsection{Bối cảnh, phạm vi (trong / ngoài)}\label{sec:se-phamvi}
\textbf{Vấn đề nghiệp vụ.} Quán cà phê quy mô nhỏ cần để khách đặt món và thanh toán không dùng tiền mặt trên điện thoại, tránh xếp hàng và nhầm đơn.

\textbf{Trong phạm vi MVP:} xem menu và chi tiết (size S/M/L, topping); giỏ hàng; đăng ký, đăng nhập, JWT, phân quyền CUSTOMER và STAFF; tạo đơn; thanh toán Ví hoặc Thẻ mô phỏng; xác nhận, theo dõi trạng thái, lịch sử đơn; máy trạng thái, thanh toán idempotent, kiểm soát tồn kho nguyên liệu đồng thời, khuyến mãi, điểm thưởng; triển khai production qua CI/CD.

\textbf{Mở rộng (nếu còn thời gian):} PWA đầy đủ (offline, thông báo); giao diện quản trị; giao diện riêng cho nhân viên quầy. \textbf{Không thực hiện:} cổng thanh toán thật, ứng dụng native, giao hàng và nhiều chi nhánh, quản lý kho nguyên liệu đầy đủ (nhập theo lô, nhà cung cấp, kiểm kê); hệ thống chỉ có kho nguyên liệu và công thức ở mức cơ bản.

\subsection{Stakeholder}\label{sec:se-stakeholder}
\begin{longtable}{@{}p{2.75cm}p{2.75cm}p{4.22cm}p{3.31cm}@{}}
\caption{Stakeholder}\label{tab:stakeholder}\\
\toprule \textbf{Stakeholder} & \textbf{Vai trò} & \textbf{Mục tiêu và kỳ vọng} & \textbf{Ảnh hưởng đến yêu cầu} \\ \midrule \endfirsthead
\toprule \textbf{Stakeholder} & \textbf{Vai trò} & \textbf{Mục tiêu và kỳ vọng} & \textbf{Ảnh hưởng đến yêu cầu} \\ \midrule \endhead
\bottomrule \endlastfoot
Khách hàng & Đặt và thanh toán đơn & Đặt nhanh, giá rõ, theo dõi được đơn & Menu, giỏ, thanh toán, lịch sử \\
Nhân viên quầy & Pha chế, cập nhật trạng thái & Nhận đúng đơn đã thanh toán & \texttt{PATCH /orders/:id/status} \\
Quản lý cửa hàng & Quản lý sản phẩm, khuyến mãi & Kiểm soát tồn kho, khuyến mãi & Tồn kho, khuyến mãi; quản trị ngoài MVP \\
Giảng viên & Khách hàng nghiệm thu & Chạy được, có tài liệu và bằng chứng & Toàn bộ sản phẩm bàn giao \\
Nhóm phát triển & Xây dựng sản phẩm & Đúng hạn, dễ bảo trì & Ràng buộc công nghệ, CI/CD \\
Đơn vị thanh toán (mô phỏng) & Xử lý Ví/Thẻ & Nhận yêu cầu hợp lệ, không trừ tiền lặp & Idempotency, trạng thái thanh toán \\
\end{longtable}

\subsection{Giả định, ràng buộc, phụ thuộc}\label{sec:se-rangbuoc}
\textbf{Giả định.} (G1) Thanh toán mặc định là mô phỏng (không có tiền thật), có thể thành công hoặc thất bại theo kịch bản; từ \texttt{v0.1.3} có thêm chế độ sandbox VNPay/MoMo khi cấu hình khóa, trong đó VNPay đã được thử đến trạng thái PAID ngày 07/11/2026. (G2) Một cửa hàng, một loại tiền (VND). (G3) Phụ phí size: S +0, M +5.000, L +10.000 đồng. (G4) Điểm thưởng: 1 điểm cho mỗi 10.000 đồng. (G5) Tồn kho tính theo nguyên liệu (\texttt{Ingredient.stock}) cho món có công thức và trừ khi thanh toán thành công; món chưa có công thức dùng \texttt{Product.stock} và trừ khi tạo đơn; hủy đơn chỉ hoàn phần đã trừ. (G6) Đơn mới ở trạng thái PENDING.

\textbf{Ràng buộc.} (R1) Bắt buộc Next.js, NestJS, PostgreSQL. (R2) Frontend mobile-first/PWA. (R3) Hạn nộp 15/11/2026. (R4) Nhóm bốn người, làm theo Scrum. (R5) Production qua HTTPS, CSDL không public, secret không nằm trong repo.

\textbf{Rủi ro ban đầu.} Không kịp triển khai production (giảm thiểu: triển khai thử sớm, 23/10); hai bài tập lớn chồng lịch (dùng chung mã nguồn); đặt hoặc thanh toán đồng thời gây âm tồn kho (cập nhật có điều kiện trong transaction, có test); sơ đồ mâu thuẫn SRS (rà bằng ma trận truy vết); thành viên chậm tiến độ (quy trình trong thỏa thuận nhóm).

\subsection{Yêu cầu chức năng}\label{sec:se-fr}
Mức ưu tiên theo MoSCoW. FR-16 (Could) và FR-17 (Won't) nằm ngoài MVP.
{\small
\begin{longtable}{@{}p{1.09cm}p{4.89cm}p{1.09cm}p{5.98cm}@{}}
\caption{Yêu cầu chức năng và tiêu chí chấp nhận}\label{tab:fr}\\
\toprule \textbf{Mã} & \textbf{Yêu cầu} & \textbf{Ưu tiên} & \textbf{Tiêu chí chấp nhận} \\ \midrule \endfirsthead
\toprule \textbf{Mã} & \textbf{Yêu cầu} & \textbf{Ưu tiên} & \textbf{Tiêu chí chấp nhận} \\ \midrule \endhead
\bottomrule \endlastfoot
FR-01 & Hệ thống hiển thị danh sách sản phẩm có tên, giá, ảnh (nếu có) & Must & \texttt{GET /\allowbreak{}products} trả danh sách; màn hình menu hiển thị đủ sản phẩm có trong seed \\
FR-02 & Hệ thống hiển thị các trạng thái loading, rỗng và lỗi của menu & Must & Mô phỏng API chậm, trả rỗng, trả lỗi thì UI hiện đúng từng trạng thái \\
FR-03 & Xem chi tiết sản phẩm; chọn size (S/M/L) và topping & Must & \texttt{GET /\allowbreak{}products/\allowbreak{}:id} trả đủ tùy chọn; giá cập nhật theo lựa chọn (S +0, M +5.000, L +10.000) \\
FR-04 & Thêm, sửa số lượng, xóa món; hiển thị tổng tiền & Must & Tổng tiền = $\Sigma$ (đơn giá + phụ phí size + topping) $\times$ số lượng \\
FR-05 & Đăng ký và đăng nhập; phát JWT; mật khẩu băm bằng bcrypt & Must & Đăng nhập đúng nhận JWT; sai mật khẩu trả 401; DB không lưu mật khẩu thô \\
FR-06 & Chặn đặt đơn khi chưa đăng nhập & Must & \texttt{POST /\allowbreak{}orders} không có JWT trả 401 \\
FR-07 & Tạo đơn từ giỏ hàng, trạng thái PENDING; kiểm tra đủ nguyên liệu cho cả đơn (món có công thức) hoặc trừ \texttt{Product.stock} trong cùng transaction (món chưa có công thức) & Must & Đơn tạo đúng; thiếu nguyên liệu hoặc hết hàng thì từ chối (400), không tạo đơn; với món có công thức, tạo đơn chưa trừ nguyên liệu \\
FR-08 & Thanh toán Ví hoặc Thẻ mô phỏng; thành công $\to$ PAID, thất bại $\to$ PAYMENT\_FAILED (API trả 200 kèm trạng thái đơn); thử lại được & Must & \texttt{POST /\allowbreak{}payments} đổi trạng thái đúng; thất bại giữ nguyên giỏ hàng \\
FR-09 & Gửi lại yêu cầu thanh toán với cùng Idempotency-Key không tạo thêm thanh toán/đơn & Must & Hai request cùng khóa cho một kết quả và một bản ghi \\
FR-10 & Máy trạng thái đơn; chặn chuyển sai trạng thái & Must & PENDING $\to$ PAID/PAYMENT\_FAILED/CANCELLED; PAID $\to$ PREPARING/CANCELLED; PREPARING $\to$ READY $\to$ COMPLETED; chuyển sai trả lỗi \\
FR-11 & Nhân viên cập nhật trạng thái đơn & Should & \texttt{PATCH /\allowbreak{}orders/\allowbreak{}:id/\allowbreak{}status} chỉ vai trò STAFF thực hiện được \\
FR-12 & Khách xem lịch sử và chi tiết đơn của mình & Must & \texttt{GET /\allowbreak{}orders/\allowbreak{}me} chỉ trả đơn của người dùng hiện tại; \texttt{GET /\allowbreak{}orders/\allowbreak{}:id} chỉ chủ đơn hoặc STAFF xem được (người khác nhận 404) \\
FR-13 & Nguyên liệu chỉ bị trừ khi thanh toán thành công (cập nhật có điều kiện trong cùng transaction); đồng thời không làm tồn kho âm; hủy đơn chỉ hoàn phần đã trừ & Must & Test N request đồng thời, số đơn thành công $\leq$ tồn kho; kho hết đúng lúc thanh toán thì không trừ tiền, không trừ kho, đơn chuyển PAYMENT\_FAILED \\
FR-14 & Áp mã khuyến mãi hợp lệ để giảm giá & Should & Mã hợp lệ giảm đúng; mã sai/hết hạn bị từ chối \\
FR-15 & Cộng điểm khi đơn thanh toán thành công (PAID): 1 điểm / 10.000đ; hủy đơn đã PAID thì trừ lại; mỗi lần đổi điểm ghi một giao dịch điểm (\texttt{LoyaltyTransaction}) & Should & \texttt{User.points} tăng đúng số điểm khi đơn PAID và giảm lại khi hủy đơn đã PAID; có bản ghi \texttt{LoyaltyTransaction} loại EARN (và REVERSAL khi hủy) tương ứng \\
FR-16 & Giao diện quản lý sản phẩm và khuyến mãi & Could & Ngoài MVP \\
FR-17 & Kết nối cổng thanh toán thật & Won't & Không thực hiện \\
\end{longtable}}

\subsection{Yêu cầu phi chức năng}\label{sec:se-nfr}
{\small
\begin{longtable}{@{}p{1.18cm}p{4.80cm}p{1.09cm}p{5.98cm}@{}}
\caption{Yêu cầu phi chức năng và cách kiểm chứng}\label{tab:nfr}\\
\toprule \textbf{Mã} & \textbf{Yêu cầu} & \textbf{Ưu tiên} & \textbf{Tiêu chí chấp nhận} \\ \midrule \endfirsthead
\toprule \textbf{Mã} & \textbf{Yêu cầu} & \textbf{Ưu tiên} & \textbf{Tiêu chí chấp nhận} \\ \midrule \endhead
\bottomrule \endlastfoot
NFR-01 & \texttt{GET /\allowbreak{}products} phản hồi nhanh với dữ liệu seed & Should & Lần chậm thứ 19 trong 20 yêu cầu \texttt{GET /\allowbreak{}products} dưới 500 ms với dữ liệu mẫu (12 sản phẩm), đo cục bộ bằng Docker Compose. Kết quả ngày 24/10/2026: 19,7 ms. Ghi chú production (Render Singapore, gói Free, đo từ Việt Nam, 30 yêu cầu): byte đầu tiên về sau 0,43 đến 0,47 giây ở 24 yêu cầu và 0,74 đến 0,78 giây ở 6 yêu cầu, trong đó khoảng 0,25 giây là độ trễ mạng một vòng \\
NFR-02 & Mật khẩu băm; JWT có hạn; kiểm tra phân quyền theo vai trò & Must & Test xác thực và phân quyền; không có mật khẩu thô trong DB hoặc log \\
NFR-03 & Secret không nằm trong repository hoặc log CI & Must & Quét repo/log không thấy secret; \texttt{.env} nằm trong \texttt{.gitignore} \\
NFR-04 & Validate dữ liệu đầu vào ở cả frontend và backend & Must & Dữ liệu sai trả lỗi 4xx rõ ràng \\
NFR-05 & Giao diện mobile-first, dùng được trên màn hình điện thoại & Must & Kiểm tra trên viewport 360–414 px, không cuộn ngang \\
NFR-06 & Thông báo lỗi dễ hiểu, không lộ chi tiết kỹ thuật & Should & Lỗi hiển thị bằng tiếng Việt, không in stack trace \\
NFR-07 & Chạy toàn bộ hệ thống bằng \texttt{docker compose up} & Must & Người ngoài nhóm chạy được theo README \\
NFR-08 & CI chạy lint, test, build cho mọi PR; pipeline CD triển khai và smoke test & Must & Có log pipeline thành công và một lần phát hiện lỗi \\
NFR-09 & \texttt{GET /\allowbreak{}health} và log để vận hành & Must & Health trả \texttt{status: ok}; smoke test gọi health sau deploy \\
NFR-10 & Migration và seed tái tạo được database & Must & Chạy migration + seed trên DB trống cho kết quả giống nhau \\
NFR-11 & Mã nguồn có lint, định dạng và test tự động & Should & CI chạy lint và test \\
\end{longtable}}
Kết quả kiểm chứng thực tế của từng NFR nằm ở mục~\ref{ch:truyvet} và mục~\ref{ch:kiemthu}; NFR-01 đo thủ công, NFR-03 mới dò sơ bộ.

\subsection{Use case và đặc tả use case}\label{sec:se-usecase}
Sơ đồ use-case tổng quát ở Hình~\ref{fig:uc-overview}; phạm vi UC-01 ở Hình~\ref{fig:uc-checkout}. UC-01 là luồng sinh doanh thu, sai ở đây gây mất tiền hoặc sai tồn kho, nên là use case trọng tâm. Hai sơ đồ theo ký hiệu chuẩn: ranh giới hệ thống là \texttt{BrewLite}, actor nối với use case bằng đường liền, tên use case chỉ nêu việc cần làm (không nêu cách làm), \texttt{<<include>>} trỏ từ use case gốc sang use case được bao gồm (Đăng nhập), \texttt{<<extend>>} trỏ từ use case mở rộng sang use case gốc và ghi điều kiện kích hoạt (khách nhập mã, thanh toán thất bại, thanh toán thành công). Hủy đơn là use case độc lập của nhân viên.
\hinh{usecase-overview.png}{.78\linewidth}{Use case tổng quan của BrewLite}{fig:uc-overview}
\nguon{nhóm vẽ bằng PlantUML (\path{docs/diagrams/usecase-overview.puml})}
\hinh{usecase-checkout.png}{.85\linewidth}{Use case của phạm vi UC-01}{fig:uc-checkout}
\nguon{nhóm vẽ bằng PlantUML (\path{docs/diagrams/usecase-checkout.puml})}

\begin{table}[htbp]\centering
\caption{Đặc tả UC-01: Đặt và thanh toán đơn hàng}\label{tab:uc01}
\begin{tabular}{@{}p{3.1cm}p{10.2cm}@{}}\toprule
Mã, tên & UC-01, Đặt và thanh toán đơn hàng \\
Actor chính / phụ & Customer / Payment Service (mô phỏng) \\
Trigger & Khách bấm ``Đặt đơn'' trong giỏ hàng \\
Yêu cầu liên quan & FR-04 đến FR-10, FR-13 đến FR-15, NFR-02, NFR-04 \\
Tiền điều kiện & Khách có JWT hợp lệ; giỏ có ít nhất một món \\
Hậu điều kiện & Thành công: đơn PAID, nguyên liệu đã trừ khi thanh toán, có bản ghi thanh toán. Thanh toán thất bại: PAYMENT\_FAILED, giữ giỏ, cho thử lại. Hết hàng hoặc lỗi: không tạo đơn, tồn kho không đổi; kho hết đúng lúc thanh toán: không trừ tiền, không trừ kho, đơn PAYMENT\_FAILED \\ \bottomrule
\end{tabular}
\end{table}
\textbf{Luồng chính.} (1) Khách mở giỏ và bấm ``Đặt đơn''. (2) Giao diện gửi \texttt{POST /orders} kèm JWT. (3) API kiểm tra JWT và validate. (4) API mở transaction: kiểm tra đủ nguyên liệu cho cả đơn và ghi nhu cầu vào \texttt{OrderIngredientUse} (món chưa có công thức thì trừ \texttt{Product.stock} có điều kiện), tạo Order (PENDING) và OrderItem, áp khuyến mãi nếu có, commit; chưa trừ nguyên liệu của món có công thức. (5) API trả mã đơn; giao diện hiển thị thanh toán. (6) Khách chọn Ví hoặc Thẻ. (7) Giao diện gửi \texttt{POST /payments} kèm \texttt{Idempotency-Key}. (8) Payment Service trả thành công. (9) Trong cùng transaction, API trừ nguyên liệu bằng cập nhật có điều kiện, ghi Payment và chuyển đơn sang PAID. (10) Giao diện hiển thị xác nhận.

\begin{longtable}{@{}p{1cm}p{5.6cm}p{6.7cm}@{}}
\caption{Luồng thay thế (A) và ngoại lệ (E) của UC-01}\label{tab:uc01-alt}\\
\toprule \textbf{Mã} & \textbf{Điều kiện} & \textbf{Xử lý} \\ \midrule \endfirsthead
\toprule \textbf{Mã} & \textbf{Điều kiện} & \textbf{Xử lý} \\ \midrule \endhead
\bottomrule \endlastfoot
A1 & Mã khuyến mãi hợp lệ & Giảm giá, cập nhật tổng tiền trước khi đặt đơn \\
A2 & Gửi lại thanh toán cùng \texttt{Idempotency-Key} & Trả kết quả cũ, không tạo thanh toán thứ hai \\
A3 & STAFF hủy đơn ở trạng thái cho phép & CANCELLED, chỉ hoàn phần nguyên liệu đã trừ, trừ lại điểm thưởng \\
E1 & Chưa đăng nhập hoặc JWT hết hạn & 401, chuyển sang đăng nhập, giữ giỏ \\
E2 & Dữ liệu không hợp lệ & 400 kèm thông báo rõ ràng \\
E3 & Món hết hàng hoặc không đủ số lượng & Rollback, không tạo đơn, báo món nào hết \\
E4 & Mã khuyến mãi sai hoặc hết hạn & Từ chối mã, cho đặt đơn không mã \\
E5 & Payment Service trả thất bại & PAYMENT\_FAILED, giữ giỏ, cho thử lại \\
E6 & Chuyển trạng thái không hợp lệ & 400, đơn giữ nguyên trạng thái \\
\end{longtable}

\section{Mô hình hóa (UML)}\label{ch:mohinh}
Nguồn PlantUML nằm ở \path{docs/diagrams}; ảnh dưới đây kết xuất từ đúng các tệp nguồn đó. Ba sơ đồ bắt buộc của đề (activity, sequence, class) cùng sơ đồ trạng thái đơn hàng được rà theo SRS và mã nguồn.

\subsection{Sơ đồ hoạt động (Activity)}\label{sec:se-activity}
Hình~\ref{fig:activity} mô tả luồng đặt và thanh toán, tách thành hai phần cho dễ đọc (Hình~\ref{fig:activity} đặt đơn, Hình~\ref{fig:activity2} thanh toán và xử lý đơn), theo bốn swimlane (Khách hàng, Hệ thống, Dịch vụ thanh toán, Nhân viên), mỗi hành động nằm đúng làn của người thực hiện. Sơ đồ có nút bắt đầu và kết thúc, nút quyết định kèm điều kiện trên từng nhánh, nút hợp nhất và hai vòng lặp: chỉnh giỏ hàng khi thiếu nguyên liệu, và thử lại khi thanh toán thất bại. Các nhánh bắt buộc có đủ: thiếu nguyên liệu, chưa đăng nhập, thanh toán thất bại, kho vừa hết lúc thanh toán. Liên quan: UC-01, FR-04 đến FR-08.
\hinhcao{activity-order-flow.png}{Sơ đồ hoạt động: đặt đơn (phần 1/2)}{fig:activity}
\hinhcao{activity-payment-flow.png}{Sơ đồ hoạt động: thanh toán và xử lý đơn (phần 2/2)}{fig:activity2}

\subsection{Sơ đồ tuần tự (Sequence)}\label{sec:se-sequence}
UC-01 được mô tả bằng hai sơ đồ tuần tự, mỗi sơ đồ một kịch bản để dễ đọc: đặt đơn (Hình~\ref{fig:seq-order}) và thanh toán (Hình~\ref{fig:seq-pay}). Lifeline xếp từ trái sang phải theo thứ tự tham gia, thanh chữ nhật mảnh trên lifeline là khoảng thời gian đối tượng đang xử lý (activation), thời gian chạy từ trên xuống dưới. Các tình huống rẽ nhánh dùng khung kết hợp: \texttt{break} cho nhánh lỗi thoát sớm, \texttt{loop} cho mỗi món trong đơn, \texttt{opt} cho bước tùy chọn và \texttt{alt} cho nhánh thay thế.

Sơ đồ đặt đơn gồm bốn lifeline: Khách hàng, Giao diện web (Next.js), API (\texttt{OrdersService}) và PostgreSQL. Sơ đồ thể hiện kiểm tra JWT và dữ liệu (401, 400), kiểm tra khóa idempotency (trả lại đơn cũ hoặc lỗi 409), kiểm tra nguyên liệu theo công thức trong một giao dịch và \texttt{ROLLBACK} khi thiếu, áp mã khuyến mãi, rồi lưu đơn ở trạng thái PENDING. Nguyên liệu chỉ được kiểm tra ở bước này; việc trừ kho diễn ra khi thanh toán thành công.

Sơ đồ thanh toán gồm năm lifeline: thêm \texttt{OrdersService} để thể hiện việc đổi trạng thái đơn. Sơ đồ thể hiện kiểm tra khóa idempotency, kiểm tra đơn hợp lệ, thử lại từ PAYMENT\_FAILED, thanh toán thành công kèm trừ nguyên liệu có điều kiện, nhánh kho vừa hết do đơn khác (\texttt{ROLLBACK}, không trừ tiền, không trừ kho, đơn chuyển PAYMENT\_FAILED) và nhánh thanh toán thất bại. Liên quan: UC-01, FR-06 đến FR-09, FR-13 đến FR-15.
\hinhcao{sequence-order.png}{Sơ đồ tuần tự: đặt đơn (UC-01)}{fig:seq-order}
\hinhcao{sequence-payment.png}{Sơ đồ tuần tự: thanh toán (UC-01)}{fig:seq-pay}

\subsection{Sơ đồ lớp / mô hình miền (Class / Domain)}\label{sec:se-class}
Hình~\ref{fig:class}, Hình~\ref{fig:class-pay} và Hình~\ref{fig:class-inv} khớp với \texttt{schema.prisma}, mỗi lớp có đủ ba ngăn (tên, thuộc tính, phương thức) với ký hiệu phạm vi truy cập (\texttt{-} riêng tư, \texttt{+} công khai). Hình~\ref{fig:class} gồm các lớp đặt hàng và điểm thưởng: User, Product, Topping, Order, OrderItem, Promotion, LoyaltyTransaction, IdempotencyRecord. Hình~\ref{fig:class-pay} gồm thanh toán, phiên QR và quà đổi điểm: Payment, PaymentIntent, PaymentAccount, Reward, Redemption (cùng User và Order ở dạng rút gọn thuộc tính). Các kiểu liệt kê (OrderStatus, PaymentMethod, PaymentStatus, PaymentProvider, IntentStatus, PromotionType, LoyaltyType, Role) không vẽ thành hộp riêng mà ghi giá trị trong ghi chú cạnh lớp dùng chúng. Quan hệ gồm kết hợp có tên, chiều đọc và multiplicity, composition giữa Order và OrderItem (dòng đơn không tồn tại khi đơn bị xóa), và lớp kết hợp \texttt{OrderItemTopping} giữa OrderItem và Topping. Hình~\ref{fig:class-inv} gồm kho nguyên liệu và công thức, với ba lớp kết hợp RecipeItem, ToppingRecipeItem và OrderIngredientUse. Các phương thức là thao tác mức miền, ánh xạ tới hàm hoặc dịch vụ trong mã (ghi chú trên sơ đồ).
\hinhngang{class-domain.png}{Sơ đồ lớp: đặt hàng và điểm thưởng}{fig:class}
\nguon{nhóm vẽ bằng PlantUML (\path{docs/diagrams/class-domain.puml})}
\hinhngang{class-payment-reward.png}{Sơ đồ lớp: thanh toán, phiên QR và quà đổi điểm}{fig:class-pay}
\nguon{nhóm vẽ bằng PlantUML (\path{docs/diagrams/class-payment-reward.puml})}
\hinhngang{class-inventory.png}{Sơ đồ lớp: kho nguyên liệu và công thức}{fig:class-inv}
\nguon{nhóm vẽ bằng PlantUML (\path{docs/diagrams/class-inventory.puml})}

\subsection{Sơ đồ trạng thái (State machine)}\label{sec:se-state}
Đơn hàng có vòng đời trạng thái (FR-10). Bảng~\ref{tab:state} mô tả các chuyển trạng thái hợp lệ do \path{apps/api/src/orders/order-state.ts} thực thi; chuyển trạng thái khác bị từ chối với mã 400 và đơn giữ nguyên.

\begin{table}[htbp]\centering
\caption{Chuyển trạng thái hợp lệ}\label{tab:state}
\begin{tabular}{@{}p{3.7cm}p{9.1cm}@{}}\toprule
\textbf{Từ} & \textbf{Sang} \\ \midrule
PENDING & PAID, PAYMENT\_FAILED, CANCELLED \\
PAYMENT\_FAILED & PAID (thử lại), CANCELLED \\
PAID & PREPARING, CANCELLED \\
PREPARING & READY \\
READY & COMPLETED \\ \bottomrule
\end{tabular}
\end{table}
Khi hủy đơn, chỉ phần nguyên liệu đã trừ được hoàn lại (đơn chưa thanh toán không đụng tới kho); nếu đơn đã PAID thì điểm thưởng đã cộng bị trừ lại (REVERSAL).

\section{Kiến trúc hệ thống}\label{ch:kientruc}
\subsection{Kiến trúc tổng thể}\label{sec:se-kientruc}
Kiến trúc phân lớp kết hợp mức container của mô hình C4: trình duyệt/PWA $\to$ frontend Next.js $\to$ backend NestJS (modular monolith, module theo domain) $\to$ PostgreSQL; dịch vụ thanh toán mô phỏng là một module trong backend (Hình~\ref{fig:arch}). Chi tiết ở \path{docs/architecture/architecture.md}.
\hinh{architecture-container.png}{.9\linewidth}{Kiến trúc mức container của BrewLite}{fig:arch}
\nguon{nhóm vẽ bằng PlantUML (\path{docs/diagrams/architecture-container.puml})}
\begin{longtable}{@{}p{2.65cm}p{2.94cm}p{4.41cm}p{3.04cm}@{}}
\caption{Thành phần (container)}\label{tab:container}\\
\toprule \textbf{Thành phần} & \textbf{Công nghệ} & \textbf{Trách nhiệm} & \textbf{Mã nguồn} \\ \midrule \endfirsthead
\toprule \textbf{Thành phần} & \textbf{Công nghệ} & \textbf{Trách nhiệm} & \textbf{Mã nguồn} \\ \midrule \endhead
\bottomrule \endlastfoot
Web & Next.js, mobile-first & Menu, chi tiết, giỏ, đăng nhập, xác nhận, lịch sử & \path{apps/web} \\
API & NestJS, JWT, validation & Nghiệp vụ, xác thực, phân quyền, truy cập CSDL & \path{apps/api} \\
CSDL & PostgreSQL 16, Prisma 7 & Người dùng, sản phẩm, đơn, thanh toán, điểm, bản ghi idempotency & \path{apps/api/prisma} \\
Payment (mô phỏng) & Module NestJS & Ví/Thẻ thành công hoặc thất bại; xử lý \texttt{Idempotency-Key} & \path{apps/api/src/payments} \\
CI/CD & GitHub Actions, GHCR & Lint, test, build, đóng gói, deploy, smoke test & \path{.github/workflows} \\
\end{longtable}
Các module backend: \texttt{auth}, \texttt{products}, \texttt{orders} (gồm khuyến mãi và điểm thưởng), \texttt{payments}, \texttt{prisma}, \texttt{health}. Vấn đề xuyên suốt: cấu hình bằng biến môi trường (\texttt{.env.example} chỉ có giá trị mẫu); log không ghi mật khẩu, token, secret; \texttt{GET /health} cho Docker healthcheck và smoke test; migration chạy trước khi API khởi động; bảo mật bằng bcrypt, JWT có hạn, validate đầu vào, CORS theo \texttt{FRONTEND\_ORIGIN}.

\textbf{Đánh đổi.} Monorepo cho phép một PR đổi cả hai phía nhưng có hai bộ \texttt{package.json}. Modular monolith dễ triển khai và kiểm thử trong một tuần nhưng khó tách dịch vụ về sau. Payment mô phỏng cho phép dựng kịch bản thành công/thất bại nhưng không phản ánh độ trễ và lỗi mạng của dịch vụ ngoài. PostgreSQL có transaction và khóa dòng phù hợp tồn kho đồng thời nhưng cần quản lý migration và sao lưu.

\subsection{Quyết định kiến trúc (ADR)}\label{sec:se-adr}
Các ADR toàn văn ở \path{docs/architecture/adr}; Bảng~\ref{tab:adr} tóm tắt.
\begin{longtable}{@{}p{1cm}p{7.6cm}p{4.7cm}@{}}
\caption{Danh sách ADR (đều ở trạng thái đã chốt; ADR-007 chọn ngày 23/10/2026)}\label{tab:adr}\\
\toprule \textbf{ADR} & \textbf{Quyết định} & \textbf{Lý do chính} \\ \midrule \endfirsthead
\toprule \textbf{ADR} & \textbf{Quyết định} & \textbf{Lý do chính} \\ \midrule \endhead
\bottomrule \endlastfoot
001 & Monorepo \path{apps/web} và \path{apps/api} & Một PR đổi cả hai phía, dùng chung CI \\
002 & PostgreSQL 16 và Prisma 7 (đã cân nhắc TypeORM) & Transaction, khóa dòng, migration \\
003 & JWT có hạn, bcrypt, phân quyền CUSTOMER/STAFF; \texttt{JWT\_SECRET} lấy từ biến môi trường & Đơn giản, dễ kiểm thử \\
004 & Payment là module NestJS; \texttt{Idempotency-Key} với ràng buộc duy nhất \texttt{(userId, key)} và \texttt{requestHash} & Không trừ tiền lặp khi gửi lại \\
005 & Trừ kho bằng \texttt{UPDATE ... WHERE stock >= n} trong transaction; nguyên liệu chỉ trừ khi thanh toán thành công & Không âm kho khi đặt và thanh toán đồng thời \\
006 & GitHub Actions và GHCR; secret trong GitHub Secrets; bật required checks \texttt{api}, \texttt{web} & Mã trên GitHub, CI chặn merge \\
007 & Render: web và API chạy image từ GHCR, PostgreSQL của Render; CD gọi Deploy Hook & HTTPS sẵn, gói miễn phí \\
\end{longtable}

\subsection{Sơ đồ triển khai (Deployment)}\label{sec:se-deployment}
Hình~\ref{fig:deploy} vẽ cấu trúc tĩnh của hệ thống theo ký hiệu chuẩn: node là hình hộp 3D có stereotype (\texttt{<<device>>}, \texttt{<<execution environment>>}, \texttt{<<database server>>}, \texttt{<<network>>}), artifact nằm trong node chứa nó, đường nối ghi giao thức, và quan hệ \texttt{<<deploy>>} từ image trên GHCR sang dịch vụ Render. Trình duyệt truy cập web và API qua Internet bằng HTTPS; API kết nối PostgreSQL của Render bằng Internal Database URL (CSDL không public). Biến môi trường và yêu cầu Public của package GHCR được ghi chú thay vì vẽ thành node. Phần quy trình CI/CD (động) được tách riêng ở Hình~\ref{fig:cicd}.
\hinh{deployment-diagram.png}{\linewidth}{Sơ đồ triển khai (deployment) của BrewLite}{fig:deploy}
\nguon{nhóm vẽ bằng PlantUML (\path{docs/architecture/deployment-diagram.puml})}
\hinhngang{cicd-pipeline.png}{Quy trình CI/CD của BrewLite (activity diagram)}{fig:cicd}
\nguon{nhóm vẽ bằng PlantUML (\path{docs/architecture/cicd-pipeline.puml}); từ \path{.github/workflows/ci.yml} và \path{cd.yml}}
Giới hạn của gói miễn phí (kiểm tra 23/10/2026): dịch vụ web ngủ khi không dùng (yêu cầu đầu tiên có thể chậm khoảng 50 giây); CSDL hết hạn ngày 22/11/2026 và không có sao lưu.

\subsection{Cơ sở dữ liệu và API}\label{sec:se-db}
Lược đồ Prisma của bản \texttt{v0.1.3} gồm 18 bảng (Hình~\ref{fig:erd}): 10 bảng lõi phục vụ luồng đặt và thanh toán theo đề (User, Product, Topping, Order, OrderItem, OrderItemTopping, Payment, IdempotencyRecord, LoyaltyTransaction, Promotion); 8 bảng còn lại thuộc chức năng bổ sung ngoài đề (nguyên liệu và công thức, thanh toán QR, đổi quà). Khuyến mãi liên kết với đơn qua mã \texttt{promoCode}, không dùng khóa ngoại.
\hinh{erd.png}{\linewidth}{Sơ đồ thực thể -- quan hệ theo \texttt{schema.prisma}}{fig:erd}
\nguon{nhóm vẽ từ \path{apps/api/prisma/schema.prisma}}
\begin{longtable}{@{}p{1.4cm}p{4cm}p{3.4cm}p{3.9cm}@{}}
\caption{Các endpoint chính}\label{tab:api}\\
\toprule \textbf{Phương thức} & \textbf{Đường dẫn} & \textbf{Quyền} & \textbf{Mã lỗi thường gặp} \\ \midrule \endfirsthead
\toprule \textbf{Phương thức} & \textbf{Đường dẫn} & \textbf{Quyền} & \textbf{Mã lỗi thường gặp} \\ \midrule \endhead
\bottomrule \endlastfoot
POST & \path{/auth/register} & Công khai & 400, 409 (trùng email) \\
POST & \path{/auth/login} & Công khai & 401 (sai thông tin) \\
GET & \path{/auth/me} & JWT & 401 \\
GET & \path{/products}, \path{/products/:id}, \path{/toppings} & Công khai & 400 (id sai), 404 \\
POST & \path{/orders} (kèm \texttt{Idempotency-Key}) & JWT & 400, 401, 409 \\
GET & \path{/orders/me}, \path{/orders/:id} & JWT (chủ đơn hoặc STAFF) & 401, 404 \\
GET & \path{/orders} & STAFF & 401, 403 \\
PATCH & \path{/orders/:id/status} & STAFF & 400 (chuyển sai), 403, 409 \\
POST & \path{/payments} (kèm \texttt{Idempotency-Key}) & JWT & 400, 401, 404, 409 \\
GET & \path{/health} & Công khai & -- \\
\end{longtable}

\section{Triển khai và tích hợp liên tục}\label{ch:trienkhai-se}
\subsection{Kho mã và quản lý tài liệu}\label{sec:se-repo}
Kho \texttt{tthan-prog/BrewLite} theo cấu trúc monorepo: \path{apps/web} (Next.js), \path{apps/api} (NestJS, Prisma), \path{tests/e2e} (Playwright), \path{docs} (yêu cầu, sơ đồ, kiến trúc, Scrum, bằng chứng, release notes), \path{.github/workflows} (\texttt{ci.yml}, \texttt{cd.yml}) và \path{docker-compose.yml}. Mỗi chức năng một nhánh; thay đổi vào \texttt{main} đi qua pull request. Nhánh \texttt{main} được bảo vệ: phải qua PR, hai check \texttt{api} và \texttt{web} phải xanh, cần một approval, không force-push. Tag \texttt{assign1-final} (commit \texttt{9581384}) đánh dấu bản Assignment 1 (mã hoàn tất 23/10/2026, thêm bản sửa ngày 07/11 để \texttt{docker compose build} chạy được trên bản clone sạch); bản phát hành \texttt{v0.1.3} có ghi chú trong \path{docs/as2/release-notes}.

\subsection{Hiện thực các chức năng chính}\label{sec:se-hienthuc}
Luồng menu $\to$ size, đường, đá $\to$ giỏ $\to$ thanh toán $\to$ xác nhận $\to$ lịch sử đã chạy trên production (23/10/2026, tài khoản \texttt{customer@brewlite.com}), kể cả thanh toán thất bại và thử lại. Các ảnh dưới đây là bằng chứng thu thập trong \path{docs/as2/evidence}.
\hinh{ev01-menu.jpg}{.92\linewidth}{Menu tải từ API trên production}{fig:menu}
\hinh{ev02-chon-size.jpg}{.92\linewidth}{Chọn size, mức đường, mức đá}{fig:size}
\hinh{ev03-gio-hang.jpg}{.92\linewidth}{Giỏ hàng: 10 ly, 34.000 đồng/ly, tạm tính 340.000 đồng}{fig:cart}
\hinh{ev04-thanh-toan.jpg}{.92\linewidth}{Thanh toán bằng thẻ mô phỏng, mã \texttt{WELCOME10}}{fig:pay}
\hinh{ev05-xac-nhan-don1-paid.jpg}{.92\linewidth}{Đơn \#1 PAID, 306.000 đồng (340.000 đồng giảm 10\%)}{fig:paid}
\hinh{ev09-thanh-toan-that-bai.jpg}{.92\linewidth}{Giả lập thanh toán thất bại: đơn \#2 PAYMENT\_FAILED, có nút ``Thử thanh toán lại''}{fig:failed}
\hinh{ev10-thu-lai-don2-paid.jpg}{.92\linewidth}{Thử lại thành công: đơn \#2 PAID, 585.000 đồng}{fig:retry}
\hinh{ev11-lich-su-2-don.jpg}{.92\linewidth}{Lịch sử có đúng hai đơn \#1 và \#2, không tạo đơn trùng khi thử lại}{fig:history}
\nguon{\path{docs/as2/evidence/01}--\path{11}, nhóm chụp ngày 23/10/2026}
Kết nối an toàn và endpoint health của production ở Hình~\ref{fig:https} và Hình~\ref{fig:health}.
\hinh{ev07-https-ket-noi-an-toan.jpg}{.82\linewidth}{Trình duyệt báo kết nối an toàn (HTTPS)}{fig:https}
\hinh{ev08-health-ok.jpg}{.82\linewidth}{\texttt{/health} trả \texttt{\{"status":"ok","service":"brewlite-backend"\}}}{fig:health}

\subsection{CI}\label{sec:se-ci}
\texttt{ci.yml} chạy cho mọi PR với bốn job: \texttt{api} (cài đặt, lint, unit test, build), \texttt{api-e2e} (migration, seed, test tích hợp trên PostgreSQL thật), \texttt{web} (lint, build) và \texttt{e2e-browser} (Playwright). Lần chạy gần nhất trên \texttt{main} (run \#85, commit \texttt{afdd005}, ngày 10/11/2026, sau khi merge PR \#58) xanh ở cả bốn job (Hình~\ref{fig:run85}).

\textbf{Minh chứng CI chặn merge (Task 4.4).} Nhóm tạo PR thử nghiệm \#30, cố ý đổi kỳ vọng một test trong \path{order-state.spec.ts} (đơn PAID không được nhảy thẳng sang COMPLETED). Kết quả: test lỗi ở máy (Hình~\ref{fig:localfail}), CI run \#13 đỏ ở job \texttt{api} (Hình~\ref{fig:run13}), PR bị chặn merge vì check \texttt{api} là Required (Hình~\ref{fig:blocked}); sau khi sửa, run \#14 xanh (Hình~\ref{fig:run14}). PR \#30 được đóng, không merge.
\hinh{ev19-test-loi-chay-o-may.jpg}{.82\linewidth}{Test lỗi khi chạy ở máy: \texttt{expected false to be true}}{fig:localfail}
\hinh{ev16-ci-do-run13-job-api-loi.jpg}{.92\linewidth}{CI run \#13 đỏ: job \texttt{api} lỗi, các job còn lại xanh}{fig:run13}
\hinh{ev17-pr30-merge-bi-chan.jpg}{.92\linewidth}{PR \#30: check \texttt{api} (Required) đỏ, ``Merging is blocked''}{fig:blocked}
\hinh{ev18-ci-xanh-run14-sau-khi-sua.jpg}{.92\linewidth}{CI run \#14 xanh sau khi sửa}{fig:run14}
\hinh{ev15-ci-xanh-merge-pr22.jpg}{.92\linewidth}{CI xanh 4/4 job (api, api-e2e, web, e2e-browser) sau khi merge PR \#22}{fig:ci4}
\hinh{ev43-ci-xanh-main-run85-pr58.jpg}{.92\linewidth}{CI run \#85 trên \texttt{main} (commit \texttt{afdd005}, 10/11/2026): xanh 4/4 job}{fig:run85}

\subsection{CD, di trú dữ liệu và quay lui}\label{sec:se-cd}
\texttt{cd.yml} chạy khi gắn tag \texttt{v*} hoặc chạy tay (\texttt{workflow\_dispatch}, tham số \texttt{image\_tag}) theo ba bước: (1) \emph{quality}: lint và test; (2) \emph{build}: dựng image Docker, đẩy lên GHCR với tag \texttt{sha-...}; (3) \emph{deploy}: gọi Deploy Hook của Render cho API và web kèm \texttt{imgURL}, rồi smoke test gọi \texttt{/health}. Secret nằm trong GitHub Secrets; migration chạy trước khi API khởi động. CD lần đầu (tag \texttt{v0.1.0}) đỏ vì thiếu cấu hình URL và hook; sau khi bổ sung, CD xanh từ \texttt{v0.1.1} và \texttt{v0.1.2}. Bản bàn giao \texttt{v0.1.3} (07/11/2026): CD chạy từ tag đỏ ở bước smoke test vì web chưa khởi động xong (\texttt{/menu} trả 502, smoke test không thử lại); nhóm chạy tay CD \#15 với \texttt{image\_tag=sha-ecd11bd} và lần này xanh (Hình~\ref{fig:cd013}).
\hinh{ev12-cd-v0.1.1-xanh.jpg}{.92\linewidth}{CD từ tag \texttt{v0.1.1}: các bước quality, build-and-push, deploy-and-smoke-test đều xanh}{fig:cd011}
\hinh{ev13-cd-v0.1.2-xanh.jpg}{.92\linewidth}{CD từ tag \texttt{v0.1.2} (commit \texttt{7b1ea96}) xanh}{fig:cd012}
\hinh{ev14-render-events-deploy-sha-7b1ea96.jpg}{.92\linewidth}{Render Events: deploy image \texttt{sha-7b1ea96} qua Deploy Hook}{fig:render}
\hinh{ev41-cd15-chay-thu-cong-image-tag-xanh.jpg}{.92\linewidth}{CD \#15 chạy tay với \texttt{image\_tag=sha-ecd11bd} (bản \texttt{v0.1.3}): xanh}{fig:cd013}

\textbf{Quay lui có minh chứng (23/10/2026).} Nhóm chạy lại workflow CD với \texttt{image\_tag} của bản ổn định trước đó (\texttt{sha-e445b63}, tức \texttt{v0.1.1}; CD \#11): bước build bị bỏ qua, chỉ deploy lại image cũ và chạy smoke test; giao diện quay về bản cũ (không có khung lỗi đỏ). Sau đó triển khai lại bản mới \texttt{sha-7b1ea96} (CD \#12), khung lỗi đỏ xuất hiện trở lại, chứng tỏ hai lần deploy khác nhau.
\hinh{ev20-rollback-cd11-bo-qua-build.jpg}{.92\linewidth}{CD \#11: bỏ qua build, deploy lại image cũ}{fig:rb1}
\hinh{ev21-rollback-cd11-log-deploy-sha-e445b63.jpg}{.92\linewidth}{Log deploy \texttt{sha-e445b63}}{fig:rb2}
\hinh{ev22-rollback-render-events-sha-e445b63.jpg}{.92\linewidth}{Render Events của lần quay lui}{fig:rb3}
\hinh{ev23-rollback-giao-dien-ban-cu-khong-co-khung-do.jpg}{.85\linewidth}{Giao diện bản cũ sau quay lui}{fig:rb4}
\hinh{ev24-deploy-lai-cd12-log-sha-7b1ea96.jpg}{.92\linewidth}{CD \#12: triển khai lại \texttt{sha-7b1ea96}}{fig:rb5}
\hinh{ev25-giao-dien-ban-moi-co-khung-do.jpg}{.92\linewidth}{Giao diện bản mới}{fig:rb6}

\section{Kiểm thử}\label{ch:kiemthu}
\subsection{Chiến lược kiểm thử}\label{sec:se-chienluoc}
Bốn mức kiểm thử, mỗi mức gắn với một nhóm yêu cầu:
\begin{itemize}
\item \textbf{Unit} (Vitest): logic thuần như máy trạng thái đơn và tính điểm thưởng.
\item \textbf{Tích hợp} (Supertest, PostgreSQL thật, job \texttt{api-e2e}): API xác thực, đơn, thanh toán, idempotency, tồn kho, khuyến mãi.
\item \textbf{E2E trình duyệt} (Playwright, job \texttt{e2e-browser}): luồng đăng ký, đặt đơn, thanh toán lỗi rồi thử lại, lịch sử; viewport di động; ba trạng thái menu; manifest PWA; giới hạn giỏ hàng.
\item \textbf{Smoke} (trong \texttt{cd.yml}): gọi \texttt{/health} sau mỗi lần deploy.
\end{itemize}
Số test lấy từ lần chạy gần nhất trên nhánh \texttt{main} (commit \texttt{2cae795}); kết quả đạt/không đạt lấy từ CI.
\begin{table}[htbp]\centering
\caption{Số lượng test theo mức}\label{tab:sotest}
\begin{tabular}{@{}p{3.6cm}p{2.3cm}p{1.5cm}p{5.4cm}@{}}\toprule
\textbf{Mức} & \textbf{Công cụ} & \textbf{Số test} & \textbf{Vị trí} \\ \midrule
Unit & Vitest & 36 & \path{apps/api/src/**/*.spec.ts} \\
Tích hợp & Supertest & 81 & \path{apps/api/test/}: \path{app.e2e-spec.ts} (2), \path{assign1.e2e-spec.ts} (30), \path{cors.e2e-spec.ts} (2), \path{qr-ingredients.e2e-spec.ts} (31), \path{realism.e2e-spec.ts} (16) \\
E2E trình duyệt & Playwright & 10 & \path{tests/e2e/brewlite-flow.spec.ts} \\ \bottomrule
\end{tabular}
\end{table}

\subsection{Test case}\label{sec:se-testcase}
Bảng~\ref{tab:testcase} liệt kê các test case chính, lấy tên từ mã nguồn kiểm thử và gắn với FR/NFR. Bảng gồm đủ 10 test Playwright (E2E-T01a đến E2E-T07).
{\small
\begin{longtable}{@{}p{1.59cm}p{6.37cm}p{1.69cm}p{3.38cm}@{}}
\caption{Test case chính và yêu cầu liên quan}\label{tab:testcase}\\
\toprule \textbf{Mã test} & \textbf{Mục tiêu (kỳ vọng)} & \textbf{Yêu cầu} & \textbf{Mức} \\ \midrule \endfirsthead
\toprule \textbf{Mã test} & \textbf{Mục tiêu (kỳ vọng)} & \textbf{Yêu cầu} & \textbf{Mức} \\ \midrule \endhead
\bottomrule \endlastfoot
INT-T01 & \texttt{GET /products} trả danh sách có id, name, price, imageUrl & FR-01 & Tích hợp \\
INT-T02 & \texttt{GET /products/:id}: có $\to$ 200, không có $\to$ 404, id sai $\to$ 400 & FR-03 & Tích hợp \\
INT-T03 & Hash bcrypt; trùng email 409, sai mật khẩu 401; đăng nhập đúng trả JWT dùng được & FR-05, NFR-02 & Tích hợp \\
INT-T04 & \texttt{POST /orders} thiếu token $\to$ 401 & FR-06 & Tích hợp \\
INT-T05 & Giá theo size; validate thiếu items, size sai, số lượng $\leq 0$ $\to$ 400; \texttt{userId} trong body bị bỏ qua & FR-07, NFR-04 & Tích hợp \\
INT-T06 & Thanh toán thành công $\to$ PAID; lỗi $\to$ PAYMENT\_FAILED, thử lại được rồi PAID & FR-08 & Tích hợp \\
INT-T07 & Hai lần cùng khóa chỉ một đơn; 10 yêu cầu song song cùng khóa; không thanh toán lần hai & FR-09 & Tích hợp \\
API-T01, INT-T08 & Mọi cặp trạng thái; chuyển sai bị chặn (400) & FR-10 & Unit, tích hợp \\
INT-T09 & STAFF đi hết PAID $\to$ PREPARING $\to$ READY $\to$ COMPLETED; CUSTOMER gọi PATCH nhận 403 & FR-11, NFR-02 & Tích hợp \\
INT-T10 & \texttt{/orders/me} chỉ trả đơn của mình; \texttt{/orders/:id}: chủ đơn 200, người khác 404 & FR-12 & Tích hợp \\
INT-T11 & 20 người đặt song song, kho 5 $\to$ đúng 5 đơn; đơn lỗi giữa chừng rollback; hủy đơn chỉ hoàn phần đã trừ & FR-13 & Tích hợp \\
INT-T12 & Mã \texttt{WELCOME10} giảm 10\% khi đủ điều kiện & FR-14 & Tích hợp \\
API-T02, INT-T13 & 1 điểm cho mỗi 10.000 đồng; PAID cộng điểm; hủy đơn đã PAID trừ lại (EARN/REVERSAL) & FR-15 & Unit, tích hợp \\
E2E-T01a & Menu hiển thị sản phẩm lấy từ API & FR-01 & E2E \\
E2E-T01b & Đăng ký, đăng nhập, thêm vào giỏ, đặt đơn, thanh toán lỗi rồi thử lại thành công, xem lịch sử & FR-04, FR-08, FR-12 & E2E \\
E2E-T02 & Viewport 390 px không cuộn ngang & NFR-05 & E2E \\
E2E-T03a/b/c & Menu đang tải, rỗng, lỗi và thử lại & FR-02, NFR-06 & E2E \\
E2E-T04 & Manifest và biểu tượng PWA & NFR-05 & E2E \\
E2E-T05 & Thanh toán QR bằng tài khoản test, hiện điểm được cộng (chức năng bổ sung ngoài đề) & FR-15 & E2E \\
E2E-T06 & Giỏ hàng không tăng quá 50 ly mỗi dòng & FR-04 & E2E \\
E2E-T07 & Ô giả lập thất bại chỉ hiện với Ví/Thẻ và gây thanh toán thất bại & FR-08 & E2E \\
\end{longtable}}

\subsection{Kết quả kiểm thử}\label{sec:se-ketquakt}
\textbf{Tự động.} Lần chạy gần nhất trên \texttt{main} (commit \texttt{2cae795}) cho 36 unit, 81 tích hợp và 10 E2E đều đạt, CI xanh ở cả bốn job. Ba test bắt buộc của Task 10 đều đạt: chặn chuyển trạng thái sai, cùng \texttt{Idempotency-Key} chỉ tạo một đơn, 20 yêu cầu song song khi kho còn 5 thì đúng 5 thành công.

\textbf{Thủ công.} Giản Thiên Lực (D) chạy trực tiếp trên backend Docker (\texttt{localhost:3001}) sáu trường hợp, tất cả đạt (Bảng~\ref{tab:manual}).
\begin{table}[htbp]\centering
\caption{Kiểm thử chức năng thủ công}\label{tab:manual}
\begin{tabular}{@{}p{6.4cm}p{6.4cm}@{}}\toprule
\textbf{Trường hợp} & \textbf{Kết quả} \\ \midrule
\texttt{GET /products} & 12 sản phẩm \\
\texttt{GET /products/1} & Cà phê sữa đá, 29.000 đồng \\
\texttt{GET /products/999999} & HTTP 404 \\
\texttt{POST /orders} không có JWT & HTTP 401 \\
\texttt{POST /orders} có JWT, thiếu \texttt{Idempotency-Key} & HTTP 400 \\
\texttt{POST /orders} hợp lệ & Đơn \#6, PENDING, 34.000 đồng \\ \bottomrule
\end{tabular}
\end{table}

\textbf{Hiệu năng (NFR-01).} Đo thủ công ngày 24/10/2026 bằng Docker Compose cục bộ với 12 sản phẩm: lần chậm thứ 19 trong 20 yêu cầu \texttt{GET /products} mất 19,7 ms, đạt ngưỡng 500 ms. Chưa có test tự động. Trên production (Render gói miễn phí, Singapore), thời gian tới byte đầu tiên đo từ Việt Nam khoảng 0,43--0,78 giây; yêu cầu đầu tiên sau thời gian không dùng có thể chậm khoảng 50 giây.

\textbf{Hạn chế đã biết.} Chưa có test tự động cho NFR-01 và NFR-03 (secret mới dò sơ bộ bằng tìm theo mẫu ngày 24/10, chưa rà thủ công đầy đủ); NFR-06 chưa có test kiểm không lộ stack trace; NFR-07 chưa có người ngoài nhóm chạy thử theo README; chưa có độ phủ (coverage) được đo.

\section{Bàn giao và đánh giá nhìn lại}\label{ch:banggiao}
\subsection{Sản phẩm bàn giao}\label{sec:se-sanpham}
\begin{longtable}{@{}p{4.2cm}p{9cm}@{}}
\caption{Danh mục bàn giao}\label{tab:sanpham}\\
\toprule \textbf{Sản phẩm} & \textbf{Vị trí / liên kết} \\ \midrule \endfirsthead
\toprule \textbf{Sản phẩm} & \textbf{Vị trí / liên kết} \\ \midrule \endhead
\bottomrule \endlastfoot
Kho mã nguồn & \url{https://github.com/tthan-prog/BrewLite} \\
Web production (HTTPS) & \url{https://brewlite-web.onrender.com} \\
API và health & \url{https://brewlite-api.onrender.com/health} \\
Bản phát hành & Tag \texttt{v0.1.3} (07/11/2026); ghi chú trong \path{docs/as2/release-notes} \\
SRS, use case, truy vết & \path{docs/requirements}, \path{docs/traceability-matrix.md} \\
Sơ đồ UML (nguồn và ảnh) & \path{docs/diagrams}, \path{docs/architecture} \\
Kiến trúc và ADR & \path{docs/architecture/architecture.md}, \path{adr/001}--\path{007} \\
CI/CD và quay lui & \path{.github/workflows/ci.yml}, \path{cd.yml}; \path{docs/deployment.md}; README mục 15 \\
Test tự động & 36 unit, 81 tích hợp, 10 E2E \\
Bằng chứng & \path{docs/as2/evidence} (ảnh 01--29 và 41) \\
Biên bản Sprint & \path{docs/as2/scrum/minutes/a2-*} \\
Video demo end-to-end & \url{https://drive.google.com/file/d/1uxrJSc9OkNkD-68vCTKXsPb0DcUudtat/view?usp=sharing} (Giản Thiên Lực thực hiện) \\
\end{longtable}
CSDL production không public (kết nối nội bộ); không có secret trong repository (\texttt{.env} bị \texttt{.gitignore} chặn, giá trị thật nằm trong GitHub Secrets và Render).

\subsection{Sprint Review / Demo}\label{sec:se-review}
Increment demo tại Review là bản đang chạy trên production lúc đó; bản bàn giao cuối là \texttt{v0.1.3} (07/11/2026). Bảng~\ref{tab:review} tóm tắt Review Sprint 2 (26/10). Sprint Goal phần lớn đạt: có URL production HTTPS, CI/CD chạy, quay lui đã chứng minh; video demo được hoàn tất sau buổi Review.
\begin{table}[htbp]\centering
\caption{Kết quả Sprint Review Sprint 2}\label{tab:review}
\begin{tabular}{@{}p{4.4cm}p{1.8cm}p{2cm}p{4.6cm}@{}}\toprule
\textbf{Hạng mục} & \textbf{Demo} & \textbf{Đạt DoD} & \textbf{Bằng chứng} \\ \midrule
5.1 Hoàn thiện MVP & Có & Có & Ảnh 01--11 \\
5.2 Tích hợp, kiểm thử & Có & Có & CI 4 job, 10 E2E \\
5.3 CD và quay lui & Có & Có & Ảnh 12--14, 20--25 \\
5.4 Production, release, README & Có & Gần đạt & Release trên GitHub; bản bàn giao \texttt{v0.1.3} \\
5.4 Video demo & Có (sau Review) & Chờ nghiệm thu & Liên kết ở Bảng~\ref{tab:sanpham} \\
5.4 Review, Retrospective & Có & Có & Biên bản \\ \bottomrule
\end{tabular}
\end{table}

\subsection{Retrospective}\label{sec:se-retro}
\textbf{Điều làm tốt, nên giữ:} quay lui bằng \texttt{image\_tag} chạy được và có minh chứng; smoke test và \texttt{/health} bắt lỗi sau triển khai; bằng chứng ảnh được đặt tên có thứ tự; mọi thay đổi vào \texttt{main} đi qua PR và CI; CI chặn merge có bằng chứng đỏ và xanh.

\textbf{Điều cần cải thiện:} CD lần đầu đỏ vì thiếu cấu hình URL và hook; Render gói miễn phí làm web ngủ và CSDL hết hạn 22/11/2026, không có sao lưu; nghiệm thu chéo chưa ghi trên PR (cột Reviews trống); một lần tạo PR nhầm nhánh đích; một lần gõ sai mã khuyến mãi (\texttt{WELCOM10}) khi thử làm tưởng lỗi giao diện.

\begin{table}[htbp]\centering
\caption{Hành động sau Sprint}\label{tab:action}
\begin{tabular}{@{}p{6.3cm}p{3.2cm}p{3.2cm}@{}}\toprule
\textbf{Hành động} & \textbf{Người làm} & \textbf{Kiểm tra vào} \\ \midrule
Quay video demo end-to-end & Giản Thiên Lực & Đã xong, liên kết ở Bảng~\ref{tab:sanpham} \\
Mở web production một lần trước giờ chấm để tránh chờ khởi động & Nguyễn Lê Nhựt Thắng & Trước buổi chấm \\
Ghi kết quả nghiệm thu chéo vào PR (comment hoặc Approve) & Người nghiệm thu theo bảng Review & Trước buổi chấm \\
Sao lưu hoặc làm mới CSDL trước 22/11 nếu còn dùng & Nguyễn Lê Nhựt Thắng & Sau khi nộp \\ \bottomrule
\end{tabular}
\end{table}

\section{Ma trận truy vết}\label{ch:truyvet}
Ma trận nối yêu cầu $\to$ use case $\to$ thiết kế $\to$ mã nguồn $\to$ test. Bản đầy đủ và được cập nhật ở \path{docs/traceability-matrix.md} (23/10/2026); cột ``Mã nguồn'' và ``Test'' trỏ tới mã và test có thật trong repo. Ô ghi ``--'' nghĩa là không áp dụng hoặc chưa làm, không điền trước. Mã test: API-Txx là unit, INT-Txx là tích hợp, E2E-Txx là Playwright. Mọi yêu cầu mức Must đều có thiết kế, mã và test, trừ các điểm ghi ở cột Trạng thái.
{\scriptsize\setlength{\tabcolsep}{3pt}\setstretch{1.0}
\begin{longtable}{@{}p{1.07cm}p{0.90cm}p{1.53cm}p{2.30cm}p{2.64cm}p{1.79cm}p{1.79cm}@{}}
\caption{Ma trận truy vết yêu cầu}\label{tab:truyvet}\\
\toprule \textbf{Yêu cầu} & \textbf{Ưu tiên} & \textbf{Use case} & \textbf{Thiết kế} & \textbf{Mã nguồn} & \textbf{Test} & \textbf{Trạng thái} \\ \midrule \endfirsthead
\toprule \textbf{Yêu cầu} & \textbf{Ưu tiên} & \textbf{Use case} & \textbf{Thiết kế} & \textbf{Mã nguồn} & \textbf{Test} & \textbf{Trạng thái} \\ \midrule \endhead
\bottomrule \endlastfoot
FR-01 & Must & US-01 & Activity, Class & \texttt{GET /products} & INT-T01, E2E-T01a & Đạt \\
FR-02 & Must & US-01 & Activity & Trang menu & E2E-T03a/b/c & Đạt \\
FR-03 & Must & US-02 & Class & \texttt{GET /products/:id} & INT-T02 & Đạt \\
FR-04 & Must & US-03 & Activity & Giỏ hàng (web) & E2E-T01b, T06 & Đạt \\
FR-05 & Must & US-04 & Class (User) & \texttt{POST /auth/*} & INT-T03 & Đạt \\
FR-06 & Must & UC-01 (E1) & Activity, Sequence & \texttt{JwtAuthGuard} & INT-T04 & Đạt \\
FR-07 & Must & UC-01 & Sequence, Class & \texttt{POST /orders} & INT-T05 & Đạt \\
FR-08 & Must & UC-01 & Sequence, Class & \texttt{POST /payments} & INT-T06, E2E-T01b, T07 & Đạt \\
FR-09 & Must & UC-01 (A2) & Sequence, ADR-004 & \texttt{Idempotency-Key} & INT-T07 & Đạt \\
FR-10 & Must & UC-01, UC-02 & Class (OrderStatus) & \path{order-state.ts} & API-T01, INT-T08 & Đạt \\
FR-11 & Should & UC-02 & Class (Role) & \texttt{PATCH /orders/:id/status} & INT-T09 & Đạt \\
FR-12 & Must & US-05 & Class & \texttt{GET /orders/me} & INT-T10 & Đạt \\
FR-13 & Must & UC-01 (E3) & Sequence, ADR-005 & Cập nhật có điều kiện & INT-T11 & Đạt \\
FR-14 & Should & UC-01 (A1) & Class (Promotion) & \texttt{promoCode} & INT-T12 & Đạt \\
FR-15 & Should & UC-02 & Class (LoyaltyTransaction) & \path{pointsForTotal} & API-T02, INT-T13 & Đạt \\
FR-16 & Could & -- & -- & -- & -- & Ngoài MVP \\
FR-17 & Won't & -- & -- & -- & -- & Không làm \\
NFR-01 & Should & -- & Architecture & \texttt{GET /products} & Đo thủ công & Đạt, chưa tự động \\
NFR-02 & Must & -- & ADR-003 & bcrypt, \path{JwtAuthGuard} & INT-T03, T09 & Đạt \\
NFR-03 & Must & -- & ADR-006 & \path{.gitignore}, GitHub Secrets & Dò sơ bộ & Chưa rà đủ \\
NFR-04 & Must & -- & -- & class-validator & INT-T05 & Đạt \\
NFR-05 & Must & -- & Kiến trúc (web) & CSS, \path{manifest.ts} & E2E-T02, T04 & Đạt, chưa service worker \\
NFR-06 & Should & -- & -- & Thông báo tiếng Việt & E2E-T03c & Một phần \\
NFR-07 & Must & -- & Deployment & \path{docker-compose.yml} & Chạy 24/10 & Chưa có người ngoài nhóm thử \\
NFR-08 & Must & -- & ADR-006 & \path{ci.yml}, \path{cd.yml} & CI \#13/\#14, CD & Đạt \\
NFR-09 & Must & -- & Deployment & \texttt{GET /health} & Smoke test & Đạt \\
NFR-10 & Must & -- & ADR-002 & \path{prisma/migrations}, seed & Job \texttt{api-e2e} & Đạt \\
NFR-11 & Should & -- & ADR-006 & Job \texttt{api}, \texttt{web} & CI & Đạt \\
\end{longtable}
}

\section{Kết luận}\label{ch:ketluan-se}
\subsection{Kết quả đạt được}\label{sec:se-dat}
Nhóm hoàn thành luồng end-to-end từ menu đến lịch sử đơn và chạy trên production qua HTTPS, gồm cả thanh toán thất bại và thử lại không tạo đơn trùng. Yêu cầu được đặc tả thành 17 FR và 11 NFR có tiêu chí chấp nhận; ba sơ đồ UML, sơ đồ trạng thái, kiến trúc, ERD, deployment và bảy ADR được rà theo mã nguồn. CI có bốn job và đã chứng minh chặn merge khi test lỗi; CD triển khai từ tag, có smoke test và quay lui đã chạy thật. Có 36 unit, 81 tích hợp và 10 E2E test; ma trận truy vết nối mọi yêu cầu Must với thiết kế, mã và test.

\subsection{Hạn chế}\label{sec:se-hanche}
\begin{itemize}
\item Nghiệm thu chéo chưa có Approve hoặc bình luận trên GitHub; PR cột Reviews còn trống.
\item NFR-01 chỉ đo thủ công một lần; NFR-03 mới dò sơ bộ; chưa đo độ phủ test.
\item Render gói miễn phí: web ngủ khi không dùng, CSDL hết hạn 22/11/2026 và không có sao lưu.
\item Thanh toán mặc định là mô phỏng; sandbox VNPay đã thử đến PAID qua Return URL (07/11/2026), còn IPN và MoMo sandbox chưa kiểm chứng; PWA mới ở mức manifest; chưa có giao diện nhân viên và quản trị.
\item Phân bố commit giữa các thành viên chưa đều.
\end{itemize}

\subsection{Hướng phát triển}\label{sec:se-huong}
Bổ sung service worker để PWA dùng được offline; giao diện barista và quản trị sản phẩm, khuyến mãi; kết nối cổng thanh toán thật; lưu mức đường và đá vào đơn; tự động hóa đo hiệu năng và quét bí mật trong CI; chuyển sang gói hạ tầng có sao lưu.

\clearpage
\appendix
\section{Bảng đối chiếu yêu cầu đề bài với báo cáo}\label{app:doichieu}
Các hạng mục lấy theo cách đánh số Task của Assignment 2 trong tài liệu của nhóm (\path{docs/as2/scrum/minutes/a2-*}).
\begin{longtable}{@{}p{5.59cm}p{4.19cm}p{3.59cm}@{}}
\toprule
\textbf{Hạng mục của đề bài} & \textbf{Mục trong báo cáo} & \textbf{Trạng thái} \\ \midrule \endfirsthead
\toprule
\textbf{Hạng mục của đề bài} & \textbf{Mục trong báo cáo} & \textbf{Trạng thái} \\ \midrule \endhead
\bottomrule \endlastfoot
1.1 Bối cảnh, stakeholder, phạm vi & mục~\ref{sec:se-phamvi}, mục~\ref{sec:se-stakeholder}, mục~\ref{sec:se-rangbuoc} & Đã đáp ứng \\
1.2 SRS (FR, NFR, ưu tiên, tiêu chí) & mục~\ref{sec:se-fr}, mục~\ref{sec:se-nfr} & Đã đáp ứng \\
1.3 Use case và đặc tả & mục~\ref{sec:se-usecase} & Đã đáp ứng \\
2 Activity, Sequence, Class & mục~\ref{sec:se-activity}, mục~\ref{sec:se-sequence}, mục~\ref{sec:se-class} & Đã đáp ứng \\
3.1 Kiến trúc và quyết định thiết kế & mục~\ref{sec:se-kientruc}, mục~\ref{sec:se-adr} & Đã đáp ứng \\
3.2 Deployment diagram & mục~\ref{sec:se-deployment} & Đã đáp ứng \\
4.1 Repository, quản lý tài liệu & mục~\ref{sec:se-repo}, mục~\ref{sec:se-dod} & Đã đáp ứng \\
4.3 MVP Sprint 1 & mục~\ref{sec:se-hienthuc}, Bảng~\ref{tab:sprint1} & Đã đáp ứng \\
4.4 CI, một lần đỏ và một lần xanh & mục~\ref{sec:se-ci} & Đã đáp ứng \\
5.1 Hoàn thiện luồng MVP & mục~\ref{sec:se-hienthuc} & Đã đáp ứng \\
5.2 Tích hợp và kiểm thử & mục~\ref{ch:kiemthu} & Đáp ứng một phần (chưa đo coverage) \\
5.3 CD và quay lui & mục~\ref{sec:se-cd} & Đã đáp ứng \\
5.4 URL, release, README, bàn giao & mục~\ref{sec:se-sanpham} & Đã đáp ứng \\
5.4 Sprint Review, Retrospective & mục~\ref{sec:se-review}, mục~\ref{sec:se-retro} & Đã đáp ứng \\
Ma trận truy vết & mục~\ref{ch:truyvet} & Đã đáp ứng \\
\end{longtable}

\section{Phân công và đóng góp}\label{app:phancong}
Tỷ lệ đóng góp lấy từ \path{docs/chung/bang-dong-gop.md}: Assignment 1 tính theo điểm ước lượng của Product Backlog (tổng 40 điểm), Assignment 2 tính theo số hạng mục trong \path{docs/as2/phan-cong.md} (tổng 33 hạng mục); tỷ lệ tổng là trung bình cộng hai tỷ lệ. Đây là khối lượng được giao, không phải số giờ làm.
\begin{table}[htbp]\centering\small
\caption{Phân công và tỷ lệ đóng góp}\label{tab:donggop}
\begin{tabular}{@{}>{\raggedright\arraybackslash}p{3.2cm}p{1.8cm}>{\raggedright\arraybackslash}p{4.0cm}p{1.05cm}p{1.05cm}p{1.05cm}@{}}\toprule
\textbf{Thành viên} & \textbf{MSSV} & \textbf{Phụ trách chính (Assignment 2)} & \textbf{AS1} & \textbf{AS2} & \textbf{Tổng} \\ \midrule
Nguyễn Lê Nhựt Thắng (A) & 3124560085 & Scrum Master; code, CI/CD, triển khai, sequence, báo cáo & 41,3\% & 42,4\% & 41,8\% \\
Hồ Thiên Ngân (B) & 3124410224 & Mô hình dữ liệu, class diagram, ERD, kiến trúc, ADR & 26,3\% & 15,2\% & 20,7\% \\
Hà Thiên Long (C) & 3124410193 & Yêu cầu, truy vết, backlog, kiểm thử, Product Owner tạm & 10,0\% & 21,2\% & 15,6\% \\
Giản Thiên Lực (D) & 3124410206 & Activity, use case, giao diện, video, slide & 22,5\% & 21,2\% & 21,9\% \\ \bottomrule
\end{tabular}
\end{table}

\section{Tư liệu bổ sung}\label{app:khac}
\subsection{Cấu trúc kho mã}\label{app:khac1}
\begin{lstlisting}[caption={Cấu trúc thư mục chính của repository},label={lst:tree}]
BrewLite/
  apps/api/        NestJS, Prisma (schema, migrations, seed), test
  apps/web/        Next.js (App Router)
  tests/e2e/       Playwright
  docs/            requirements, diagrams, architecture (ADR),
                   traceability-matrix, deployment, as1, as2, chung
  .github/workflows/  ci.yml, cd.yml
  docker-compose.yml
\end{lstlisting}

\subsection{Liên kết}\label{app:khac2}
\begin{itemize}\raggedright\setlength{\itemsep}{0pt}
\item Kho mã: \url{https://github.com/tthan-prog/BrewLite}
\item Web production: \url{https://brewlite-web.onrender.com}
\item API: \url{https://brewlite-api.onrender.com/health}
\item Video demo: \url{https://drive.google.com/file/d/1uxrJSc9OkNkD-68vCTKXsPb0DcUudtat/view?usp=sharing}
\end{itemize}

\clearpage
\begin{thebibliography}{80}

\bibitem{scrum}
Ken Schwaber, Jeff Sutherland,
``\textbf{The Scrum Guide}'',
\textit{link: https://scrumguides.org/scrum-guide.html}, 2020, lần truy cập cuối: 26/10/2026.

\bibitem{de2}
Khoa Công nghệ thông tin, Trường Đại học Sài Gòn,
``\textbf{BrewLite -- Assignment 2: Requirements to Deployment (VER 1.0)}'',
\textit{môn Công nghệ phần mềm, học kỳ I năm học 2026--2027}.

\end{thebibliography}
\end{document}
